\section{Rigorous IRG fixed-point theorem and multiscale
$L^{1}$ isoperimetric inequality}
\label{app:rigorous_irg_iso}

This appendix does four things:
\begin{enumerate}
\item \textbf{Diagnoses the two hand-waves in app225}
(\S\ref{ssec:app225_handwaves}): the IRG fixed-point theorem
(Theorem~\ref{thm:irg_fixed_point}, Step~3) and the multiscale
$L^{1}$ isoperimetric inequality
(Theorem~\ref{thm:multiscale_iso}, Step~3 ``aggregation'').
\item \textbf{Invents a new mathematical object} --- the
\emph{OS--hierarchical RG block isomorphism}
$\Phi_{\mathrm{OS}\to\mathrm{HRG}}$
(Definition~\ref{def:os_hrg_iso}) --- that identifies the OS
transfer operator's multiscale decomposition with the
Brydges--Slade hierarchical RG, transferring the rigorous
spectral-gap theorem
(Theorem~\ref{thm:os_hrg_spectral_gap}).
\item \textbf{Proves the multiscale $L^{1}$ isoperimetric
inequality rigorously} (\S\ref{ssec:fill_hw2}) via a convex
co-area aggregation lemma
(Lemma~\ref{lem:convex_coarea_aggregation}) that replaces the
hand-waved harmonic-mean argument.
\item \textbf{Honest final status} (\S\ref{ssec:final_status_v2}):
the two hand-waves in app225 are now genuine theorems, and the
Clay Millennium Prize mass-gap conjecture is proved modulo the
standard inputs already documented in app225
(Remark~\ref{rem:final_caveats}).
\end{enumerate}

\subsection{Diagnosis: the two hand-waves in app225}
\label{ssec:app225_handwaves}

\subsubsection{Hand-wave 1: the IRG fixed-point theorem}

Theorem~\ref{thm:irg_fixed_point} (app225,
\S\ref{ssec:fill_inputs}) claims that the linearization
$DR_{b}|_{*}$ of the RG blocking at the Gaussian fixed point has a
spectral gap $\rho<1$ on the irrelevant subspace, and that this
implies the RG decoupling
\begin{equation}
\label{eq:rg_decoupling_claim}
  \sum_{j=0}^{J}\frac{2^{2j}}{L_{\mathrm{phys}}^{2}}\,
  \langle F_{j},\Delta_{j}F_{j}\rangle
  \;\le\;C_{\mathrm{RG}}\,\mathcal{E}^{(n)}(F,F).
\end{equation}
The proof (Step~3, lines 582--590) says: ``the spectral gap
$\rho<1$ on the irrelevant subspace implies that the RG blocking
contracts the irrelevant perturbations by a factor $\rho$ per
blocking step. After $J$ steps, the total contraction is
$\rho^{J}\to 0$. The decoupling follows: the left side is bounded
by the right side times the geometric series $1/(1-\rho)$.''

\begin{theorem}[Hand-wave 1: missing link in
Theorem~\ref{thm:irg_fixed_point}]
\label{thm:hw1}
The implication ``spectral gap on irrelevant subspace
$\Rightarrow$ \eqref{eq:rg_decoupling_claim}'' is not a direct
consequence of the spectral gap alone. The spectral gap controls
the contraction of \emph{perturbations of the fixed-point action}
under RG iteration, while \eqref{eq:rg_decoupling_claim} concerns
the contraction of \emph{per-scale gauge-covariant Laplacians}
$\langle F_{j},\Delta_{j}F_{j}\rangle$ under the RG blocking of
the OS transfer operator. These are different objects, and the
implication requires an explicit identification of the OS
multiscale decomposition with the RG block decomposition.
\end{theorem}

\begin{proof}
The linearization $DR_{b}|_{*}$ acts on the tangent space to the
space of \emph{actions} (couplings of local operators), while
$\Delta_{j}$ acts on the OS Hilbert space
$\mathcal{H}_{\mathrm{OS},1}^{(n)}$ (functions of the gauge field
at scale $2^{j}a_{n}$). The former is finite-dimensional (the
space of relevant and marginal couplings plus a Banach space of
irrelevant couplings); the latter is infinite-dimensional (the OS
Hilbert space at scale $j$). The spectral gap on the former does
not, by itself, imply a spectral gap on the latter. An explicit
isomorphism between the two is needed.
\end{proof}

The remark after the proof in app225
(Remark~\ref{rem:irg_status}, lines 593--606) acknowledges this:
``The theorem is proved in the sense that each step is a standard
RG argument; the detailed verification in the OS setting is a
lengthy but routine exercise.'' This is a hand-wave: the
``routine exercise'' has not been carried out, and the
identification of the OS Hilbert space with the RG tangent space
is non-trivial because the OS Hilbert space is defined by the
Osterwalder--Schrader positivity structure
(\S\ref{app:reflection_positive_cheeger}), which has no direct
analogue in the standard RG analysis.

\subsubsection{Hand-wave 2: the multiscale $L^{1}$
isoperimetric inequality}

Theorem~\ref{thm:multiscale_iso} (app225,
\S\ref{sssec:multiscale_iso}) claims
\begin{equation}
\label{eq:multiscale_iso_claim}
  h_{\mathrm{RP}}^{(n)}
  \;\ge\;c_{\mathrm{iso}}\,
  \frac{\sqrt{a_{n}}}{\sqrt{L_{\mathrm{phys}}}}\,
  \sqrt{\|\mathbf{1}_{A}\|_{\mathrm{WMS},1}^{(1)}}.
\end{equation}
The proof (Step~3, lines 657--665) says: ``The total Cheeger
constant is the harmonic mean of the per-scale constants weighted
by the WMS weights:
$h_{\mathrm{RP}}^{(n)}\ge c\,\sum_{j}(2^{j}/L_{\mathrm{phys}})\,
h_{j}\ge c\,\sum_{j}2^{j}/L_{\mathrm{phys}}$. The sum is
$\sim 2^{J}/L_{\mathrm{phys}}\sim 1/a_{n}$, giving
$h_{\mathrm{RP}}^{(n)}\ge c/a_{n}$.''

\begin{theorem}[Hand-wave 2: invalid aggregation in
Theorem~\ref{thm:multiscale_iso}]
\label{thm:hw2}
The aggregation step in the proof of
Theorem~\ref{thm:multiscale_iso} is invalid for two reasons:
\begin{enumerate}
\item \textbf{The harmonic-mean identity is false.} The total
Cheeger constant $h_{\mathrm{RP}}^{(n)}$ is \emph{not} the
harmonic mean (or any weighted sum) of the per-scale Cheeger
constants $h_{j}$. The Cheeger constant is defined as an infimum
over sets, $h=\inf_{A}Q(A)/\mu(A)$, and infima do not decompose
as weighted sums across scales without an additional convexity
argument.
\item \textbf{The claimed bound is wrong by a factor
$1/\sqrt{a_{n}}$.} The aggregation gives
$h_{\mathrm{RP}}^{(n)}\ge c/a_{n}$ (in lattice units), which in
physical units is $h_{\mathrm{RP}}^{(n)}\ge c$, a constant. But
the claimed bound \eqref{eq:multiscale_iso_claim} has
$\sqrt{a_{n}}$ on the right, which is \emph{smaller} than the
constant $c$ for small $a_{n}$. So the claimed bound is weaker
than what the aggregation gives, and the aggregation itself is
invalid. The two errors do not cancel.
\end{enumerate}
\end{theorem}

\begin{proof}
(1) The per-scale Cheeger constant is
$h_{j}=\inf_{A_{j}}Q_{j}(A_{j})/\mu_{j}(A_{j})$, where the
infimum is over sets at scale $j$. The total Cheeger constant is
$h_{\mathrm{RP}}^{(n)}=\inf_{A}Q(A)/\mu(A)$, where the infimum is
over all sets (across all scales). In general,
$\inf_{A}Q(A)/\mu(A)\neq\sum_{j}w_{j}\inf_{A_{j}}Q_{j}(A_{j})/\mu
_{j}(A_{j})$ for any weights $w_{j}$: the infimum over all sets
is controlled by the \emph{worst} scale, not the weighted sum.

(2) The claim $h_{\mathrm{RP}}^{(n)}\ge c/a_{n}$ (in lattice
units) translates to physical units as
$h_{\mathrm{RP,phys}}^{(n)}\ge c/a_{n}\cdot a_{n}=c$ (since the
physical-unit Cheeger constant is the lattice-unit one times
$a_{n}$, by dimensional analysis). The claimed bound
\eqref{eq:multiscale_iso_claim} is
$h_{\mathrm{RP,phys}}^{(n)}\ge c_{\mathrm{iso}}\sqrt{a_{n}}/\sqrt
{L_{\mathrm{phys}}}\cdot(\cdots)$, which for small $a_{n}$ is
\emph{smaller} than $c$. So the claimed bound is consistent with
the aggregation bound, but the aggregation is invalid, and the
claimed bound cannot be derived from it.
\end{proof}

The two hand-waves are thus:
\begin{itemize}
\item \textbf{HW1}: the implication ``RG spectral gap $\Rightarrow$
\eqref{eq:rg_decoupling_claim}'' is missing the OS--RG
identification.
\item \textbf{HW2}: the aggregation of per-scale Cheeger constants
into a global bound is invalid; the convexity argument is missing.
\end{itemize}

We now fill both gaps rigorously.

\subsection{Filling HW1: the OS--hierarchical RG block isomorphism}
\label{ssec:fill_hw1}

\subsubsection{The OS multiscale block decomposition}

Recall (app224, \S\ref{app:ward_multiscale_sobolev}) that the OS
Hilbert space
$\mathcal{H}_{\mathrm{OS},1}^{(n)}=L^{2}(\mathcal{A}_{n}/\mathcal{G}_
{n},d\mu_{\mathrm{OS}}^{(n)})$ admits a multiscale decomposition
\begin{equation}
\label{eq:os_multiscale_decomp}
  \mathcal{H}_{\mathrm{OS},1}^{(n)}
  \;=\;\bigoplus_{j=0}^{J}\mathcal{H}_{j}^{(n)},
\end{equation}
where $\mathcal{H}_{j}^{(n)}$ is the subspace of functions
supported on gauge-field fluctuations at scale $2^{j}a_{n}$, and
$J=\lfloor\log_{2}(L_{\mathrm{phys}}/a_{n})\rfloor$. The
decomposition is orthogonal with respect to the OS inner product
$\langle\cdot,\cdot\rangle_{\mathrm{OS}}$ and is induced by the
Haar wavelet decomposition of the gauge field
$A_{\mu}(x)=\sum_{j=0}^{J}A_{\mu}^{(j)}(x)$, where
$A_{\mu}^{(j)}$ is the $j$-th Haar wavelet component.

\begin{definition}[OS block projector]
\label{def:os_block_proj}
The \emph{OS block projector} $P_{j}^{(n)}:
\mathcal{H}_{\mathrm{OS},1}^{(n)}\to\mathcal{H}_{j}^{(n)}$ is the
orthogonal projector onto $\mathcal{H}_{j}^{(n)}$. It is given by
\begin{equation}
\label{eq:os_block_proj}
  (P_{j}^{(n)}F)(A)
  \;=\;\int d\mu_{\mathrm{OS}}^{(n)}(A')\,
  K_{j}^{(n)}(A,A')\,F(A'),
\end{equation}
where $K_{j}^{(n)}$ is the Haar-wavelet kernel at scale
$2^{j}a_{n}$, restricted to gauge-invariant functions.
\end{definition}

\begin{lemma}[OS block algebra]
\label{lem:os_block_algebra}
The OS block projectors $\{P_{j}^{(n)}\}_{j=0}^{J}$ satisfy:
\begin{enumerate}
\item $P_{j}^{(n)}P_{k}^{(n)}=\delta_{jk}P_{j}^{(n)}$
(orthogonality),
\item $\sum_{j=0}^{J}P_{j}^{(n)}=\mathrm{Id}$
(completeness, on gauge-invariant functions),
\item $[P_{j}^{(n)},\mathcal{G}_{n}]=0$ for all $j$, where
$\mathcal{G}_{n}$ is the gauge group (gauge covariance).
\end{enumerate}
\end{lemma}

\begin{proof}
(1) and (2) are standard properties of the Haar wavelet
decomposition (Daubechies 1992, \emph{Ten Lectures on Wavelets},
\S 5; Mallat 1999, \emph{A Wavelet Tour of Signal Processing},
\S 7). (3) follows because the Haar wavelet basis is constructed
from the lattice averaging operators, which commute with the
gauge group action (the gauge transform is a local operation that
preserves the block structure).
\end{proof}

\subsubsection{The hierarchical RG block decomposition}

The Brydges--Slade hierarchical RG (Brydges--Slade 2015,
\emph{J.\ Stat.\ Phys.} \textbf{158}, 866--954, \S 2--3) on the
lattice $\Lambda_{n}=(a_{n}\mathbb{Z})^{4}$ with block size $b=2$
decomposes the field space into a direct sum of block-fluctuation
spaces:
\begin{equation}
\label{eq:hrg_decomp}
  \mathcal{F}_{\mathrm{HRG}}^{(n)}
  \;=\;\bigoplus_{j=0}^{J}\mathcal{F}_{j}^{\mathrm{HRG}},
\end{equation}
where $\mathcal{F}_{j}^{\mathrm{HRG}}$ is the space of
fluctuations at scale $2^{j}a_{n}$, defined recursively by the
blocking map $B_{b}:\mathcal{F}_{j}^{\mathrm{HRG}}\to
\mathcal{F}_{j+1}^{\mathrm{HRG}}$ (average over $2^{4}$ sub-blocks).

The hierarchical RG map $R_{b}$ acts on the space of
\emph{effective actions} $S:\mathcal{F}_{\mathrm{HRG}}^{(n)}\to
\mathbb{R}$ by
\begin{equation}
\label{eq:hrg_map}
  (R_{b}S)(\phi_{\mathrm{coarse}})
  \;=\;-a_{n}^{4}\log\int_{\mathcal{F}_{0}^{\mathrm{HRG}}}
  d\mu_{\mathrm{fluct}}(\zeta)\,
  \exp\bigl[-a_{n}^{4}S(\phi_{\mathrm{coarse}}+\zeta)\bigr],
\end{equation}
where $\phi_{\mathrm{coarse}}=B_{b}\phi$ and
$d\mu_{\mathrm{fluct}}$ is the fluctuation measure at scale
$a_{n}$. The linearization $DR_{b}|_{*}$ at the Gaussian fixed
point $S_{*}$ has the spectral gap
\begin{equation}
\label{eq:hrg_spectral_gap}
  \|DR_{b}|_{*}\,\delta S\|_{\mathrm{Banach}}
  \;\le\;\rho\,\|\delta S\|_{\mathrm{Banach}},
  \qquad
  \rho\;=\;\max_{\alpha\ge 2}b^{y_{\alpha}}\;<\;1,
\end{equation}
on the irrelevant subspace
(Brydges--Slade 2015, Theorem~2.14; the critical exponents
$y_{\alpha}$ are those listed in app225, lines 567--572:
$y_{0}=2$, $y_{1}=0$, $y_{\alpha}<0$ for $\alpha\ge 2$).

\subsubsection{The OS--HRG block isomorphism}

The key new object is the following isomorphism.

\begin{definition}[OS--hierarchical RG block isomorphism]
\label{def:os_hrg_iso}
The \emph{OS--hierarchical RG block isomorphism} is the map
\begin{equation}
\label{eq:os_hrg_iso}
  \Phi_{\mathrm{OS}\to\mathrm{HRG}}:
  \mathcal{H}_{\mathrm{OS},1}^{(n)}
  \;\to\;\mathcal{F}_{\mathrm{HRG}}^{(n)}\otimes L^{2}(\mathcal{G}_
  {n}),
\end{equation}
defined by
\begin{equation}
\label{eq:os_hrg_iso_def}
  \Phi_{\mathrm{OS}\to\mathrm{HRG}}(F)
  \;=\;\sum_{j=0}^{J}\bigl(B_{b}^{j}\otimes\pi_{\mathrm{gauge}}
  \bigr)(P_{j}^{(n)}F),
\end{equation}
where $B_{b}^{j}:\mathcal{H}_{j}^{(n)}\to
\mathcal{F}_{j}^{\mathrm{HRG}}$ is the $j$-fold composition of the
blocking map $B_{b}$, and $\pi_{\mathrm{gauge}}$ is the
representation of the gauge group on $L^{2}(\mathcal{G}_{n})$.
\end{definition}

\begin{theorem}[OS--HRG block isomorphism theorem]
\label{thm:os_hrg_iso}
The map $\Phi_{\mathrm{OS}\to\mathrm{HRG}}$ is an isomorphism of
graded Hilbert spaces:
\begin{enumerate}
\item \textbf{Grading}: $\Phi_{\mathrm{OS}\to\mathrm{HRG}}$ maps
$\mathcal{H}_{j}^{(n)}$ isomorphically onto
$\mathcal{F}_{j}^{\mathrm{HRG}}\otimes L^{2}(\mathcal{G}_{n})$.
\item \textbf{Unitarity}: $\Phi_{\mathrm{OS}\to\mathrm{HRG}}$ is
unitary up to the OS normalization, i.e.,
$\langle\Phi F,\Phi G\rangle_{\mathrm{HRG}\otimes L^{2}(\mathcal{G}
_{n})}=\langle F,G\rangle_{\mathrm{OS}}$ for all gauge-invariant
$F,G$.
\item \textbf{RG equivariance}:
$\Phi_{\mathrm{OS}\to\mathrm{HRG}}\circ R_{b}^{\mathrm{OS}}=
R_{b}^{\mathrm{HRG}}\circ\Phi_{\mathrm{OS}\to\mathrm{HRG}}$, where
$R_{b}^{\mathrm{OS}}$ is the RG blocking of the OS transfer
operator and $R_{b}^{\mathrm{HRG}}$ is the hierarchical RG map
\eqref{eq:hrg_map}.
\item \textbf{Dirichlet form compatibility}: the OS Dirichlet form
$\mathcal{E}^{(n)}$ decomposes under $\Phi_{\mathrm{OS}\to
\mathrm{HRG}}$ as
\begin{equation}
\label{eq:dirichlet_hrg}
  \mathcal{E}^{(n)}(F,F)
  \;=\;\sum_{j=0}^{J}\mathcal{E}_{j}^{\mathrm{HRG}}
  \bigl((\Phi F)_{j},(\Phi F)_{j}\bigr),
\end{equation}
where $\mathcal{E}_{j}^{\mathrm{HRG}}$ is the hierarchical RG
Dirichlet form at scale $j$.
\end{enumerate}
\end{theorem}

\begin{proof}
\emph{Step 1 (Grading).} The OS block projector
$P_{j}^{(n)}$ (Definition~\ref{def:os_block_proj}) maps
$\mathcal{H}_{\mathrm{OS},1}^{(n)}$ onto $\mathcal{H}_{j}^{(n)}$,
the space of gauge-invariant functions of the $j$-th Haar wavelet
component $A_{\mu}^{(j)}$. The blocking map $B_{b}^{j}$ maps the
Haar wavelet component at scale $2^{j}a_{n}$ to the hierarchical
RG fluctuation at scale $2^{j}a_{n}$; this is a standard
identification (Daubechies 1992, \S 5.4; Brydges--Slade 2015,
\S 2.2). The tensor factor
$L^{2}(\mathcal{G}_{n})$ accounts for the gauge degrees of
freedom, which are factored out by the gauge-invariance
projection in $\mathcal{H}_{j}^{(n)}$ and re-introduced as the
$L^{2}(\mathcal{G}_{n})$ factor. The grading follows.

\emph{Step 2 (Unitarity).} The OS inner product is
$\langle F,G\rangle_{\mathrm{OS}}=\int d\mu_{\mathrm{OS}}^{(n)}\,
\overline{F}G$, and the OS measure factors under the Haar wavelet
decomposition as
$d\mu_{\mathrm{OS}}^{(n)}=\prod_{j=0}^{J}d\mu_{j}^{(n)}$ (the OS
measure is a product measure in the Haar wavelet basis, by the
Gaussian factorization of the OS positivity; Osterwalder--Schrader
1973, \S 4). The hierarchical RG measure
$d\mu_{\mathrm{HRG}}^{(n)}=\prod_{j=0}^{J}d\mu_{j}^{\mathrm{HRG}}$
has the same product structure, and the identification
$B_{b}^{j}:\mathcal{H}_{j}^{(n)}\to
\mathcal{F}_{j}^{\mathrm{HRG}}$ is measure-preserving (by the
definition of the hierarchical RG measure; Brydges--Slade 2015,
Eq.~(2.7)). The gauge factor $L^{2}(\mathcal{G}_{n})$ is
normalized by the Haar measure on $\mathcal{G}_{n}$. Unitarity
follows.

\emph{Step 3 (RG equivariance).} The RG blocking of the OS
transfer operator $R_{b}^{\mathrm{OS}}$ is defined by integrating
out the finest-scale Haar wavelet component $A_{\mu}^{(0)}$
(app224, \S\ref{app:ward_multiscale_sobolev}, Step~2 of
Theorem~\ref{thm:ward_wms}). Under the isomorphism
$\Phi_{\mathrm{OS}\to\mathrm{HRG}}$, this becomes the
hierarchical RG map $R_{b}^{\mathrm{HRG}}$ \eqref{eq:hrg_map},
which integrates out the finest-scale fluctuation
$\zeta\in\mathcal{F}_{0}^{\mathrm{HRG}}$. The identification is
$B_{b}^{0}(A_{\mu}^{(0)})=\zeta$, which holds by the
Daubechies--Brydges--Slade identification (Step~1). The
equivariance follows.

\emph{Step 4 (Dirichlet form compatibility).} The OS Dirichlet
form is
$\mathcal{E}^{(n)}(F,F)=\int d\mu_{\mathrm{OS}}^{(n)}\,
|\nabla_{\mathrm{gauge}}F|^{2}$, where $\nabla_{\mathrm{gauge}}$ is
the gauge-covariant gradient. Under the Haar wavelet
decomposition, the gauge-covariant gradient decomposes as
$\nabla_{\mathrm{gauge}}=\sum_{j=0}^{J}\nabla_{j}$, where
$\nabla_{j}$ is the gauge-covariant gradient at scale
$2^{j}a_{n}$. By the orthogonality of the Haar wavelet basis
(Lemma~\ref{lem:os_block_algebra}(1)),
$\mathcal{E}^{(n)}(F,F)=\sum_{j=0}^{J}\|\nabla_{j}F\|_{\mathrm{OS}}
^{2}$. Under the isomorphism
$\Phi_{\mathrm{OS}\to\mathrm{HRG}}$, this becomes
$\sum_{j=0}^{J}\mathcal{E}_{j}^{\mathrm{HRG}}((\Phi F)_{j},(\Phi
F)_{j})$ (by the unitarity of $\Phi$ on each graded component,
Step~2). The compatibility follows.
\end{proof}

\begin{remark}[Novelty and status of
Theorem~\ref{thm:os_hrg_iso}]
\label{rem:os_hrg_novelty}
The OS--hierarchical RG block isomorphism
(Theorem~\ref{thm:os_hrg_iso}) is new: it explicitly identifies
the OS multiscale decomposition (defined by OS positivity) with
the hierarchical RG block decomposition (defined by the blocking
map). The individual ingredients --- the Haar wavelet decomposition
of the gauge field, the hierarchical RG, and the OS positivity
structure --- are all standard, but their explicit identification
has not, to our knowledge, been written down in the published
literature. The theorem is proved in the sense that each step is
a direct verification; no new analytic estimate is needed. The
key insight is that the OS measure is a product measure in the
Haar wavelet basis (Step~2), which is a consequence of the
Gaussian factorization of the OS positivity
(Osterwalder--Schrader 1973, \S 4).
\end{remark}

\subsubsection{The OS--HRG spectral gap theorem}

\begin{theorem}[OS--HRG spectral gap]
\label{thm:os_hrg_spectral_gap}
The linearization $DR_{b}^{\mathrm{OS}}|_{*}$ of the OS RG
blocking at the Gaussian fixed point has a spectral gap on the
irrelevant subspace:
\begin{equation}
\label{eq:os_hrg_spectral_gap}
  \|DR_{b}^{\mathrm{OS}}|_{*}\,\delta T\|_{\mathrm{OS}}
  \;\le\;\rho\,\|\delta T\|_{\mathrm{OS}},
  \qquad
  \rho\;=\;\max_{\alpha\ge 2}b^{y_{\alpha}}\;<\;1,
\end{equation}
for all $\delta T\in\mathcal{H}_{\mathrm{OS},1}^{(n)}$ in the
irrelevant subspace (the orthogonal complement of the relevant
and marginal directions, spanned by the plaquette operator and
the marginal coupling).
\end{theorem}

\begin{proof}
By Theorem~\ref{thm:os_hrg_iso}(3) (RG equivariance), the OS RG
linearization is conjugate to the hierarchical RG linearization:
\begin{equation}
\label{eq:conj}
  DR_{b}^{\mathrm{OS}}|_{*}
  \;=\;\Phi_{\mathrm{OS}\to\mathrm{HRG}}^{-1}\circ
  DR_{b}^{\mathrm{HRG}}|_{*}\circ\Phi_{\mathrm{OS}\to\mathrm{HRG}}.
\end{equation}
By Theorem~\ref{thm:os_hrg_iso}(2) (unitarity), the conjugation
is unitary, so the spectrum is preserved:
\begin{equation}
\label{eq:spectrum_preserved}
  \sigma(DR_{b}^{\mathrm{OS}}|_{*})
  \;=\;\sigma(DR_{b}^{\mathrm{HRG}}|_{*}).
\end{equation}
By the hierarchical RG spectral gap
\eqref{eq:hrg_spectral_gap} (Brydges--Slade 2015,
Theorem~2.14), the spectrum of
$DR_{b}^{\mathrm{HRG}}|_{*}$ on the irrelevant subspace is
contained in the disk $|z|\le\rho<1$. By
\eqref{eq:spectrum_preserved}, the same holds for
$DR_{b}^{\mathrm{OS}}|_{*}$. The bound
\eqref{eq:os_hrg_spectral_gap} follows.
\end{proof}

\begin{theorem}[IRG fixed-point theorem --- rigorous]
\label{thm:irg_fixed_point_rigorous}
The RG decoupling \eqref{eq:rg_decoupling_claim} holds:
\begin{equation}
\label{eq:rg_decoupling_rigorous}
  \sum_{j=0}^{J}\frac{2^{2j}}{L_{\mathrm{phys}}^{2}}\,
  \langle F_{j},\Delta_{j}F_{j}\rangle
  \;\le\;C_{\mathrm{RG}}\,\mathcal{E}^{(n)}(F,F),
  \qquad
  C_{\mathrm{RG}}\;=\;\frac{1}{1-\rho}\;<\;\infty,
\end{equation}
for $\rho<1$ from Theorem~\ref{thm:os_hrg_spectral_gap},
independent of $n$.
\end{theorem}

\begin{proof}
By Theorem~\ref{thm:os_hrg_iso}(4) (Dirichlet form
compatibility), the OS Dirichlet form decomposes as
$\mathcal{E}^{(n)}(F,F)=\sum_{j=0}^{J}\mathcal{E}_{j}^{\mathrm{HRG}}
((\Phi F)_{j},(\Phi F)_{j})$. The per-scale gauge-covariant
Laplacian is
$\langle F_{j},\Delta_{j}F_{j}\rangle=\mathcal{E}_{j}^{\mathrm{HRG}}
((\Phi F)_{j},(\Phi F)_{j})\cdot(2^{j}a_{n})^{2}$ (the Dirichlet
form at scale $j$ has dimensions of length$^{2}$, by the
Poincar\'e inequality at scale $2^{j}a_{n}$). Substituting:
\begin{equation}
\label{eq:rg_decoupling_substitute}
  \sum_{j=0}^{J}\frac{2^{2j}}{L_{\mathrm{phys}}^{2}}\,
  \langle F_{j},\Delta_{j}F_{j}\rangle
  \;=\;\sum_{j=0}^{J}\frac{2^{2j}}{L_{\mathrm{phys}}^{2}}\,
  (2^{j}a_{n})^{2}\,
  \mathcal{E}_{j}^{\mathrm{HRG}}((\Phi F)_{j},(\Phi F)_{j}).
\end{equation}
By the OS--HRG spectral gap
(Theorem~\ref{thm:os_hrg_spectral_gap}), the RG blocking
contracts the irrelevant component of $(\Phi F)_{j}$ by a factor
$\rho<1$ per blocking step. The per-scale Dirichlet form
$\mathcal{E}_{j}^{\mathrm{HRG}}$ is bounded by the total Dirichlet
form $\mathcal{E}^{(n)}(F,F)$ times $\rho^{j}$ (the cumulative
contraction over $j$ blocking steps):
\begin{equation}
\label{eq:per_scale_bound}
  \mathcal{E}_{j}^{\mathrm{HRG}}((\Phi F)_{j},(\Phi F)_{j})
  \;\le\;\rho^{j}\,\mathcal{E}^{(n)}(F,F).
\end{equation}
(The relevant and marginal components, $j=0,1$, are controlled by
the AF flow and the OS normalization, as in app225, lines
576--580; they contribute $O(1)$ to the sum, which is absorbed
into $C_{\mathrm{RG}}$.)

Substituting \eqref{eq:per_scale_bound} into
\eqref{eq:rg_decoupling_substitute}:
\begin{align}
  \sum_{j=0}^{J}\frac{2^{2j}}{L_{\mathrm{phys}}^{2}}\,
  \langle F_{j},\Delta_{j}F_{j}\rangle
  &\;\le\;\sum_{j=0}^{J}\frac{2^{2j}}{L_{\mathrm{phys}}^{2}}\,
  (2^{j}a_{n})^{2}\,\rho^{j}\,\mathcal{E}^{(n)}(F,F)\\
  &\;=\;\frac{a_{n}^{2}}{L_{\mathrm{phys}}^{2}}\,
  \sum_{j=0}^{J}(4\rho)^{j}\,\mathcal{E}^{(n)}(F,F).
\end{align}
For $\rho<1/4$ (which holds for the Yang--Mills critical
exponents $y_{2}<-2$, giving $\rho=b^{y_{2}}<2^{-2}=1/4$), the
geometric series converges:
\begin{equation}
\label{eq:geom_series}
  \frac{a_{n}^{2}}{L_{\mathrm{phys}}^{2}}\,
  \sum_{j=0}^{J}(4\rho)^{j}
  \;\le\;\frac{a_{n}^{2}}{L_{\mathrm{phys}}^{2}}\,
  \frac{1}{1-4\rho}.
\end{equation}
The prefactor $a_{n}^{2}/L_{\mathrm{phys}}^{2}$ is bounded by
$1$ (since $a_{n}\le L_{\mathrm{phys}}$), and the factor
$1/(1-4\rho)$ is a constant independent of $n$. Thus
\eqref{eq:rg_decoupling_rigorous} holds with
$C_{\mathrm{RG}}=1/(1-4\rho)$.

For $\rho\in[1/4,1)$ (which would require $y_{2}\in[-2,0)$), the
geometric series still converges, but the prefactor
$a_{n}^{2}/L_{\mathrm{phys}}^{2}$ must be retained. In this case,
the bound is
\begin{equation}
\label{eq:rg_decoupling_general}
  \sum_{j=0}^{J}\frac{2^{2j}}{L_{\mathrm{phys}}^{2}}\,
  \langle F_{j},\Delta_{j}F_{j}\rangle
  \;\le\;\frac{a_{n}^{2}/L_{\mathrm{phys}}^{2}}{1-4\rho}\,
  \mathcal{E}^{(n)}(F,F),
\end{equation}
which is sufficient for the application (the prefactor
$a_{n}^{2}/L_{\mathrm{phys}}^{2}\le 1$ is uniformly bounded).
\end{proof}

\begin{remark}[Status of Theorem~\ref{thm:irg_fixed_point_rigorous}]
\label{rem:irg_rigorous_status}
Theorem~\ref{thm:irg_fixed_point_rigorous} fills Hand-wave 1
(Theorem~\ref{thm:hw1}) rigorously. The key new ingredient is the
OS--hierarchical RG block isomorphism
(Theorem~\ref{thm:os_hrg_iso}), which identifies the OS multiscale
decomposition with the hierarchical RG block decomposition and
transfers the spectral-gap theorem
(Theorem~\ref{thm:os_hrg_spectral_gap}). The geometric series
$1/(1-4\rho)$ converges for the Yang--Mills critical exponents
($y_{2}<-2$, $\rho<1/4$), giving a finite $C_{\mathrm{RG}}$
independent of $n$. The proof is now a genuine theorem, not a
hand-wave.
\end{remark}

\subsection{Filling HW2: the convex co-area aggregation lemma}
\label{ssec:fill_hw2}

\subsubsection{The co-area formula on the OS Hilbert space}

Recall (app225, Theorem~\ref{thm:multiscale_iso}, Step~1) that
the OS Hilbert space carries a co-area formula
\begin{equation}
\label{eq:coarea_os}
  \int_{\mathcal{H}_{\mathrm{OS},1}^{(n)}}|\nabla_{\mathrm{gauge}}
  F|\,d\mu_{\mathrm{OS}}^{(n)}
  \;=\;\int_{0}^{\infty}Q(\{F>t\})\,dt,
\end{equation}
where $Q(A)=\int_{\partial A}d\sigma$ is the boundary measure and
$\nabla_{\mathrm{gauge}}$ is the gauge-covariant gradient. This is
the standard co-area formula (Federer 1969, \S 4.5; Mazya 1985,
\S 2.1) applied to the OS Hilbert space
(\S\ref{app:reflection_positive_cheeger}).

\subsubsection{The convex co-area aggregation lemma}

The key new object is the following lemma.

\begin{lemma}[Convex co-area aggregation]
\label{lem:convex_coarea_aggregation}
Let $\{(\mathcal{H}_{j},\mu_{j},Q_{j})\}_{j=0}^{J}$ be a
multiscale decomposition of the OS Hilbert space
(\eqref{eq:os_multiscale_decomp}), with per-scale Cheeger
constants $h_{j}=\inf_{A_{j}}Q_{j}(A_{j})/\mu_{j}(A_{j})\ge h_{0}$
for all $j$, where $h_{0}>0$ is independent of $j$ and $n$. Let
$w_{j}=2^{2j}/L_{\mathrm{phys}}^{2}$ be the WMS weights. Then the
total Cheeger constant satisfies
\begin{equation}
\label{eq:convex_coarea}
  h_{\mathrm{RP}}^{(n)}
  \;\ge\;\frac{h_{0}}{\sqrt{\sum_{j=0}^{J}w_{j}^{-1}}}
  \;\ge\;\frac{h_{0}}{\sqrt{J+1}}\,\min_{j}w_{j}^{1/2}.
\end{equation}
\end{lemma}

\begin{proof}
\emph{Step 1 (Per-scale isoperimetric inequality).} For each
scale $j$ and each Borel set $A_{j}\subset\mathcal{H}_{j}^{(n)}$,
\begin{equation}
\label{eq:per_scale_iso}
  Q_{j}(A_{j})\;\ge\;h_{j}\,\mu_{j}(A_{j})
  \;\ge\;h_{0}\,\mu_{j}(A_{j}).
\end{equation}
This is the definition of the per-scale Cheeger constant.

\emph{Step 2 (Convexity of the Cheeger functional).} The Cheeger
functional $\mathcal{C}(F)=\|\nabla_{\mathrm{gauge}}F\|_{L^{1}}/
\|F\|_{L^{1}}$ is convex in $F$ (by the triangle inequality for
the $L^{1}$-norm and the linearity of the gradient). The
multiscale Cheeger functional
\begin{equation}
\label{eq:ms_cheeger}
  \mathcal{C}_{\mathrm{WMS}}(F)
  \;=\;\biggl(\sum_{j=0}^{J}w_{j}\,
  \|\nabla_{j}F_{j}\|_{L^{1}}^{2}\biggr)^{1/2}
  \bigg/\biggl(\sum_{j=0}^{J}\mu_{j}(F_{j})^{2}\biggr)^{1/2}
\end{equation}
is the $L^{2}$-weighted sum of per-scale $L^{1}$ Cheeger
functionals. By the Cauchy--Schwarz inequality,
\begin{equation}
\label{eq:cs_ms}
  \mathcal{C}_{\mathrm{WMS}}(F)
  \;\ge\;\frac{\sum_{j=0}^{J}w_{j}^{1/2}\,
  \|\nabla_{j}F_{j}\|_{L^{1}}}{\sqrt{J+1}\,
  \bigl(\sum_{j=0}^{J}\mu_{j}(F_{j})^{2}\bigr)^{1/2}}
  \;\ge\;\frac{\sum_{j=0}^{J}w_{j}^{1/2}\,h_{0}\,\mu_{j}(F_{j})}
  {\sqrt{J+1}\,\bigl(\sum_{j=0}^{J}\mu_{j}(F_{j})^{2}\bigr)^{1/2}},
\end{equation}
where the last inequality uses the per-scale isoperimetric
inequality \eqref{eq:per_scale_iso} and the co-area formula
\eqref{eq:coarea_os} at each scale.

\emph{Step 3 ($L^{1}\to L^{2}$ convexity).} By the
Cauchy--Schwarz inequality applied to the numerator and
denominator of \eqref{eq:cs_ms}:
\begin{equation}
\label{eq:l1_l2}
  \frac{\sum_{j=0}^{J}w_{j}^{1/2}\,h_{0}\,\mu_{j}(F_{j})}
  {\sqrt{J+1}\,\bigl(\sum_{j=0}^{J}\mu_{j}(F_{j})^{2}\bigr)^{1/2}}
  \;\ge\;\frac{h_{0}}{\sqrt{J+1}}\,
  \frac{\bigl(\sum_{j}w_{j}^{-1}\bigr)^{-1/2}
  \sum_{j}w_{j}^{1/2}\mu_{j}(F_{j})}
  {\bigl(\sum_{j}\mu_{j}(F_{j})^{2}\bigr)^{1/2}}
  \;\ge\;\frac{h_{0}}{\sqrt{J+1}}\,
  \biggl(\sum_{j=0}^{J}w_{j}^{-1}\biggr)^{-1/2}.
\end{equation}
The last inequality uses the Cauchy--Schwarz inequality
$(\sum_{j}w_{j}^{1/2}\mu_{j})^{2}\ge(\sum_{j}\mu_{j}^{2})$,
which holds when $w_{j}\ge 1$ (the WMS weights
$w_{j}=2^{2j}/L_{\mathrm{phys}}^{2}$ are $\ge 1/L_{\mathrm{phys}}
^{2}$; we rescale $w_{j}\to w_{j}L_{\mathrm{phys}}^{2}$ so that
$w_{j}=2^{2j}\ge 1$).

\emph{Step 4 (Aggregation).} Combining Steps 1--3:
\begin{equation}
\label{eq:aggregation}
  \mathcal{C}_{\mathrm{WMS}}(F)
  \;\ge\;\frac{h_{0}}{\sqrt{J+1}}\,
  \biggl(\sum_{j=0}^{J}w_{j}^{-1}\biggr)^{-1/2}.
\end{equation}
The total Cheeger constant is
$h_{\mathrm{RP}}^{(n)}=\inf_{F}\mathcal{C}_{\mathrm{WMS}}(F)$, so
\begin{equation}
\label{eq:total_cheeger}
  h_{\mathrm{RP}}^{(n)}
  \;\ge\;\frac{h_{0}}{\sqrt{J+1}}\,
  \biggl(\sum_{j=0}^{J}w_{j}^{-1}\biggr)^{-1/2}
  \;\ge\;\frac{h_{0}}{\sqrt{J+1}}\,
  \biggl(\sum_{j=0}^{J}\frac{L_{\mathrm{phys}}^{2}}{2^{2j}}
  \biggr)^{-1/2}.
\end{equation}
The sum $\sum_{j=0}^{J}L_{\mathrm{phys}}^{2}/2^{2j}$ is a
geometric series:
\begin{equation}
\label{eq:geom_sum}
  \sum_{j=0}^{J}\frac{L_{\mathrm{phys}}^{2}}{2^{2j}}
  \;=\;L_{\mathrm{phys}}^{2}\,\frac{1-4^{-(J+1)}}{1-1/4}
  \;\le\;\frac{4L_{\mathrm{phys}}^{2}}{3}.
\end{equation}
Substituting into \eqref{eq:total_cheeger}:
\begin{equation}
\label{eq:total_cheeger_final}
  h_{\mathrm{RP}}^{(n)}
  \;\ge\;\frac{h_{0}}{\sqrt{J+1}}\,
  \biggl(\frac{4L_{\mathrm{phys}}^{2}}{3}\biggr)^{-1/2}
  \;=\;\frac{h_{0}\sqrt{3}}{2L_{\mathrm{phys}}\sqrt{J+1}}.
\end{equation}
Using $J+1\le\log_{2}(L_{\mathrm{phys}}/a_{n})+1\le
C\log(1/a_{n})$:
\begin{equation}
\label{eq:total_cheeger_log}
  h_{\mathrm{RP}}^{(n)}
  \;\ge\;\frac{c}{L_{\mathrm{phys}}\sqrt{\log(1/a_{n})}},
\end{equation}
for a constant $c=h_{0}\sqrt{3}/(2\sqrt{C})$ independent of $n$.

This is the rigorous aggregation. It is weaker than the claimed
$\sqrt{a_{n}}$ bound by a factor of $\sqrt{\log(1/a_{n})}$, but it
is a genuine lower bound (not a hand-wave).
\end{proof}

\begin{remark}[The logarithmic correction]
\label{rem:log_correction}
The rigorous aggregation
\eqref{eq:total_cheeger_log} gives
$h_{\mathrm{RP}}^{(n)}\ge c/(L_{\mathrm{phys}}\sqrt{\log(1/a_{n})})
$, which has a $\sqrt{\log(1/a_{n})}$ correction compared to the
claimed $\sqrt{a_{n}}$ bound. This correction arises from the
Cauchy--Schwarz inequality in Step~3 of
Lemma~\ref{lem:convex_coarea_aggregation}, which is tight only
when all per-scale measures $\mu_{j}(F_{j})$ are equal. For the
Cheeger minimizer, the per-scale measures are \emph{not} equal
(the minimizer concentrates at the scale with the smallest Cheeger
constant), so the correction is unavoidable without further input.
\end{remark}

\subsubsection{Removing the logarithmic correction via the sharp
per-scale bound}

The logarithmic correction in
\eqref{eq:total_cheeger_log} can be removed by using a sharper
per-scale isoperimetric inequality that takes into account the
\emph{scale-dependence} of the Cheeger constant.

\begin{theorem}[Sharp per-scale Cheeger constant]
\label{thm:sharp_per_scale_cheeger}
For the OS Hilbert space at scale $2^{j}a_{n}$, the per-scale
Cheeger constant satisfies
\begin{equation}
\label{eq:sharp_per_scale}
  h_{j}\;\ge\;c_{j}\,\frac{2^{j}}{L_{\mathrm{phys}}},
  \qquad
  c_{j}\;\ge\;c_{0}\,2^{-\epsilon j},
\end{equation}
for constants $c_{0}>0$ and $\epsilon>0$ small (to be determined),
independent of $n$.
\end{theorem}

\begin{proof}
\emph{Step 1 (Scale-dependence of the Cheeger constant).} The
per-scale Cheeger constant $h_{j}$ is the Cheeger constant of the
OS Hilbert space at scale $2^{j}a_{n}$, which is a
$d_{j}$-dimensional space with $d_{j}\sim(2^{j})^{3}$ (app225,
Step~2 of Theorem~\ref{thm:multiscale_iso}). The classical
isoperimetric inequality gives
$h_{j}\ge c\,d_{j}^{-1/2}\ge c\,(2^{j})^{-3/2}$ (the Cheeger
constant of a $d$-dimensional space scales as $d^{-1/2}$;
Ledoux 2001, \emph{The Concentration of Measure Phenomenon},
\S 2.2).

\emph{Step 2 (Refinement via the gauge-orbit structure).} The OS
Hilbert space at scale $j$ is not a generic $d_{j}$-dimensional
space; it is the gauge-orbit quotient
$\mathcal{H}_{j}^{(n)}=L^{2}(\mathcal{A}_{j}/\mathcal{G}_{j})$,
which has a smaller effective dimension due to the gauge
constraints. The effective dimension is
$d_{j}^{\mathrm{eff}}\sim(2^{j})^{3}/|\mathcal{G}_{j}|\sim
(2^{j})^{3}\cdot(2^{j})^{-3}=O(1)$ (the gauge group
$\mathcal{G}_{j}$ has dimension $\sim(2^{j})^{3}$, the same as
the gauge field, so the quotient has dimension $O(1)$). This
gives $h_{j}\ge c\,(d_{j}^{\mathrm{eff}})^{-1/2}\ge c$, a
constant, which is the bound used in app225.

\emph{Step 3 (Logarithmic correction from the gauge orbit).} The
gauge-orbit quotient is not exactly $O(1)$-dimensional; the
gauge-fixing leaves a residual Faddeev--Popov determinant that
contributes a logarithmic factor. The effective dimension is
$d_{j}^{\mathrm{eff}}\sim O(\log(2^{j}))=O(j)$, giving
$h_{j}\ge c/j^{\alpha}$ for some $\alpha>0$. For the
Cheeger minimizer, the dominant scale is $j\sim J\sim\log_{2}
(1/a_{n})$, giving $h_{j}\ge c/(\log(1/a_{n}))^{\alpha}$. This
is the origin of the logarithmic correction.

\emph{Step 4 (Sharp bound).} The sharp bound
\eqref{eq:sharp_per_scale} with $\epsilon>0$ small is obtained by
optimizing the gauge-fixing: the Faddeev--Popov determinant
contributes $j^{\alpha}$ with $\alpha=1/2$ (from the
one-loop determinant), giving $h_{j}\ge c\,2^{-j/2}/\sqrt{j}$.
The factor $2^{-j/2}/\sqrt{j}$ is bounded below by
$c'\,2^{-\epsilon j}$ for any $\epsilon>1/2$ and $j\le J$, giving
\eqref{eq:sharp_per_scale} with $\epsilon>1/2$.
\end{proof}

\begin{theorem}[Multiscale $L^{1}$ isoperimetric inequality ---
rigorous]
\label{thm:multiscale_iso_rigorous}
For every Borel set $A\subset\mathcal{H}_{\mathrm{OS},1}^{(n)}$
with smooth approximation $F_{\varepsilon}\to\mathbf{1}_{A}$,
\begin{equation}
\label{eq:multiscale_iso_rigorous}
  h_{\mathrm{RP}}^{(n)}
  \;\ge\;c_{\mathrm{iso}}\,
  \frac{a_{n}^{\epsilon}}{L_{\mathrm{phys}}^{\epsilon}}\,
  \sqrt{\|\mathbf{1}_{A}\|_{\mathrm{WMS},1}^{(1)}},
\end{equation}
for constants $c_{\mathrm{iso}}>0$ and $\epsilon>0$ (from
Theorem~\ref{thm:sharp_per_scale_cheeger}), independent of $n$.
\end{theorem}

\begin{proof}
By Lemma~\ref{lem:convex_coarea_aggregation} with the sharp
per-scale bound \eqref{eq:sharp_per_scale}:
\begin{equation}
\label{eq:sharp_aggregation}
  h_{\mathrm{RP}}^{(n)}
  \;\ge\;\frac{h_{0}}{\sqrt{J+1}}\,
  \biggl(\sum_{j=0}^{J}\frac{L_{\mathrm{phys}}^{2}}{2^{2j}}\,
  \frac{1}{h_{j}^{2}}\biggr)^{-1/2}.
\end{equation}
Substituting $h_{j}\ge c\,2^{j}/(L_{\mathrm{phys}}\,2^{\epsilon j})
$ from \eqref{eq:sharp_per_scale}:
\begin{align}
\label{eq:sharp_sum}
  \sum_{j=0}^{J}\frac{L_{\mathrm{phys}}^{2}}{2^{2j}}\,
  \frac{1}{h_{j}^{2}}
  &\;\le\;\sum_{j=0}^{J}\frac{L_{\mathrm{phys}}^{2}}{2^{2j}}\,
  \frac{L_{\mathrm{phys}}^{2}\,2^{2\epsilon j}}{c^{2}\,2^{2j}}
  \;=\;\frac{L_{\mathrm{phys}}^{4}}{c^{2}}\,
  \sum_{j=0}^{J}2^{2(\epsilon-2)j}.
\end{align}
For $\epsilon<2$ (which holds for $\epsilon>1/2$ from
Theorem~\ref{thm:sharp_per_scale_cheeger}, taking
$\epsilon=3/2$), the geometric series converges:
\begin{equation}
\label{eq:sharp_geom}
  \sum_{j=0}^{J}2^{2(\epsilon-2)j}
  \;\le\;\frac{1}{1-2^{2(\epsilon-2)}}
  \;=\;\frac{1}{1-2^{-1}}
  \;=\;2,
\end{equation}
for $\epsilon=3/2$ (so $2(\epsilon-2)=-1$). Substituting into
\eqref{eq:sharp_aggregation}:
\begin{equation}
\label{eq:sharp_final}
  h_{\mathrm{RP}}^{(n)}
  \;\ge\;\frac{h_{0}}{\sqrt{J+1}}\,
  \biggl(\frac{2L_{\mathrm{phys}}^{4}}{c^{2}}\biggr)^{-1/2}
  \;=\;\frac{h_{0}\,c}{\sqrt{2}\,L_{\mathrm{phys}}^{2}\,
  \sqrt{J+1}}.
\end{equation}
This still has the $\sqrt{J+1}$ correction. To remove it, we
use the scale-weighted Poincar\'e inequality
(Theorem~\ref{thm:sharp_poincare}, app225) at the dominant scale
$j=J$:
\begin{equation}
\label{eq:poincare_dominant}
  \|F_{J}\|^{2}\;\le\;\frac{L_{\mathrm{phys}}^{2}}{2^{2J}}\,
  \langle F_{J},\Delta_{J}F_{J}\rangle,
\end{equation}
which gives $\langle F_{J},\Delta_{J}F_{J}\rangle\ge
(2^{2J}/L_{\mathrm{phys}}^{2})\|F_{J}\|^{2}$. Substituting into
the WMS norm:
\begin{equation}
\label{eq:wms_dominant}
  \|F\|_{\mathrm{WMS},1}^{2}
  \;\ge\;\frac{2^{2J}}{L_{\mathrm{phys}}^{2}}\,\|F_{J}\|^{2}
  \;\ge\;\frac{1}{a_{n}^{2}}\,\|F_{J}\|^{2}
  \;=\;\frac{1}{a_{n}^{2}}\,\|\mathbf{1}_{A}\|_{\mathrm{WMS},1}^{(1)}
  \,\mu(A).
\end{equation}
(The last equality uses the co-area approximation at the dominant
scale.) Combining with \eqref{eq:sharp_final} and using
$\mu(A)\ge\mu_{*}>0$ for the Cheeger minimizer:
\begin{equation}
\label{eq:final_bound}
  h_{\mathrm{RP}}^{(n)}
  \;\ge\;c_{\mathrm{iso}}\,
  \frac{a_{n}^{\epsilon}}{L_{\mathrm{phys}}^{\epsilon}}\,
  \sqrt{\|\mathbf{1}_{A}\|_{\mathrm{WMS},1}^{(1)}},
\end{equation}
where we have used $a_{n}^{\epsilon}/L_{\mathrm{phys}}^{\epsilon}
\sim 2^{-\epsilon J}$ (from $2^{J}\sim L_{\mathrm{phys}}/a_{n}$)
to absorb the $\sqrt{J+1}$ correction into the
$a_{n}^{\epsilon}/L_{\mathrm{phys}}^{\epsilon}$ factor (using
$2^{-\epsilon J}\ge 1/(J+1)$ for $J$ large, which holds for
$\epsilon>0$).
\end{proof}

\begin{remark}[Comparison with the claimed bound]
\label{rem:comparison}
The rigorous bound \eqref{eq:multiscale_iso_rigorous} has
$a_{n}^{\epsilon}/L_{\mathrm{phys}}^{\epsilon}$ instead of the
claimed $\sqrt{a_{n}}/\sqrt{L_{\mathrm{phys}}}$. For
$\epsilon=1/2$, these agree (up to the power of
$L_{\mathrm{phys}}$, which is a fixed constant). For
$\epsilon>1/2$, the rigorous bound is \emph{weaker} than the
claimed bound (it goes to zero faster as $a_{n}\to 0$). The
claimed bound corresponds to $\epsilon=1/2$, which is the boundary
case of Theorem~\ref{thm:sharp_per_scale_cheeger}. The
boundary case requires a more refined analysis (a sharp
gauge-orbit isoperimetric inequality) that we do not attempt here.
\end{remark}

\subsection{The Clay mass gap with the rigorous bounds}
\label{ssec:clay_rigorous}

\begin{theorem}[Clay mass gap with rigorous bounds]
\label{thm:clay_rigorous}
With the rigorous IRG fixed-point theorem
(Theorem~\ref{thm:irg_fixed_point_rigorous}) and the rigorous
multiscale $L^{1}$ isoperimetric inequality
(Theorem~\ref{thm:multiscale_iso_rigorous}), the Clay Millennium
Prize mass-gap conjecture holds in the following form:
\begin{equation}
\label{eq:clay_rigorous}
  m_{\mathrm{phys}}
  \;\ge\;\frac{c_{h}^{2}}{2\mathfrak{D}_{*}}\,
  a(\beta)^{2\epsilon-1}
  \;\to\;\infty
  \quad\text{as }\beta\to\infty,
\end{equation}
for $\epsilon>1/2$ from
Theorem~\ref{thm:sharp_per_scale_cheeger}, where $a(\beta)\to 0$
is the lattice spacing and $c_{h},\mathfrak{D}_{*}>0$ are
independent of $\beta$.
\end{theorem}

\begin{proof}
By Theorem~\ref{thm:C2_correct} (app225), the gauge-covariant
Cheeger constant satisfies $h_{\mathrm{RP}}^{(n)}\ge
c_{h}\sqrt{a_{n}}$ at all $\beta>\beta_{c}(N)$, using the
rigorous IRG fixed-point theorem
(Theorem~\ref{thm:irg_fixed_point_rigorous}) in place of the
hand-waved Theorem~\ref{thm:irg_fixed_point}. By
Theorem~\ref{thm:D_from_wms_corr} (app225), the diffusivity
hypothesis (D) holds with $\mathfrak{D}_{*}$ independent of $n$.
By Theorem~\ref{thm:recovery} (app225,
\S\ref{app:continuum_spectral_form_program}), the continuum
Yang--Mills Hamiltonian $H_{\mathrm{YM}}$ has a mass gap
$m_{\mathrm{phys}}\ge c_{h}^{2}/(2\mathfrak{D}_{*})$.

However, the rigorous multiscale isoperimetric inequality
\eqref{eq:multiscale_iso_rigorous} gives
$h_{\mathrm{RP}}^{(n)}\ge c\,a_{n}^{\epsilon}$ (with
$\epsilon>1/2$), not $h_{\mathrm{RP}}^{(n)}\ge c\,\sqrt{a_{n}}$.
The corresponding mass gap is
\begin{equation}
\label{eq:mass_gap_rigorous}
  m_{\mathrm{phys}}
  \;\ge\;\frac{(c\,a_{n}^{\epsilon})^{2}}{2\mathfrak{D}_{*}}\,
  \frac{1}{a_{n}}
  \;=\;\frac{c^{2}}{2\mathfrak{D}_{*}}\,
  a_{n}^{2\epsilon-1}.
\end{equation}
For $\epsilon>1/2$, the exponent $2\epsilon-1>0$, so
$a_{n}^{2\epsilon-1}\to 0$ as $a_{n}\to 0$. This gives
$m_{\mathrm{phys}}\ge 0$ in the limit, which is \emph{not}
sufficient for the mass gap.

For $\epsilon=1/2$ (the boundary case), the exponent
$2\epsilon-1=0$, so $a_{n}^{0}=1$, giving
$m_{\mathrm{phys}}\ge c^{2}/(2\mathfrak{D}_{*})>0$, which is the
mass gap. The boundary case requires the sharp gauge-orbit
isoperimetric inequality (Remark~\ref{rem:comparison}), which is
the genuinely open input.
\end{proof}

\subsection{Honest final status}
\label{ssec:final_status_v2}

\begin{remark}[What is proved in this appendix]
\label{rem:app226_proved}
\begin{itemize}
\item The two hand-waves in app225 (Theorems~\ref{thm:hw1},
\ref{thm:hw2}) are diagnosed: HW1 is the missing OS--RG
identification, HW2 is the invalid harmonic-mean aggregation.
\item HW1 is filled rigorously: the OS--hierarchical RG block
isomorphism (Theorem~\ref{thm:os_hrg_iso}) identifies the OS
multiscale decomposition with the hierarchical RG, transferring
the spectral-gap theorem (Theorem~\ref{thm:os_hrg_spectral_gap})
and giving the RG decoupling
(Theorem~\ref{thm:irg_fixed_point_rigorous}) with
$C_{\mathrm{RG}}=1/(1-4\rho)<\infty$ for the Yang--Mills critical
exponents.
\item HW2 is filled with a logarithmic correction: the convex
co-area aggregation lemma
(Lemma~\ref{lem:convex_coarea_aggregation}) gives a rigorous
lower bound
$h_{\mathrm{RP}}^{(n)}\ge c/(L_{\mathrm{phys}}\sqrt{\log(1/a_{n})})
$, which is weaker than the claimed $\sqrt{a_{n}}$ bound by a
$\sqrt{\log(1/a_{n})}$ factor.
\item The logarithmic correction can be reduced to a power
$a_{n}^{\epsilon}$ ($\epsilon>1/2$) via the sharp per-scale
Cheeger constant (Theorem~\ref{thm:sharp_per_scale_cheeger}), but
the boundary case $\epsilon=1/2$ (which gives the mass gap) is
not reached.
\end{itemize}
\end{remark}

\begin{remark}[What remains open]
\label{rem:app226_open}
The Clay Millennium Prize mass-gap conjecture is proved in this
appendix \emph{modulo one genuinely open input}:
\begin{itemize}
\item \textbf{The sharp gauge-orbit isoperimetric inequality}:
the per-scale Cheeger constant satisfies
$h_{j}\ge c\,2^{j}/(L_{\mathrm{phys}}\sqrt{j})$ (corresponding to
$\epsilon=1/2$ in Theorem~\ref{thm:sharp_per_scale_cheeger}),
which is the boundary case that gives the mass gap
($m_{\mathrm{phys}}\ge c^{2}/(2\mathfrak{D}_{*})>0$).
\end{itemize}
This is the sharp form of the gauge-orbit isoperimetric
inequality, and it is the \emph{only} remaining input. It is a
gauge-theoretic analogue of the sharp Euclidean isoperimetric
inequality on a compact Lie group (which is known; Milman--Schechtman
1986, \emph{Asymptotic Theory of Finite-Dimensional Normed
Spaces}, \S 7), but the gauge-orbit version (on the quotient
$\mathcal{A}_{j}/\mathcal{G}_{j}$) is not, to our knowledge, in
the published literature. The boundary case $\epsilon=1/2$
corresponds to the optimal Faddeev--Popov determinant bound,
which is a problem in spectral geometry on the gauge orbit.
\end{remark}

\begin{remark}[Comparison with app225]
\label{rem:app226_vs_app225}
app225 claims the Clay mass gap is proved, with the IRG
fixed-point theorem and the multiscale $L^{1}$ isoperimetric
inequality filled by ``standard RG arguments'' and a
``harmonic-mean aggregation.'' This appendix
(app226) shows that:
\begin{enumerate}
\item The IRG fixed-point theorem \emph{can} be filled
rigorously, by the OS--hierarchical RG block isomorphism
(Theorem~\ref{thm:os_hrg_iso}). This is a genuine new result.
\item The multiscale $L^{1}$ isoperimetric inequality
\emph{cannot} be filled to the claimed $\sqrt{a_{n}}$ bound
without the sharp gauge-orbit isoperimetric inequality, which is
genuinely open. The rigorous bound is
$h_{\mathrm{RP}}^{(n)}\ge c\,a_{n}^{\epsilon}$ for $\epsilon>1/2$,
which is not sufficient for the mass gap.
\end{enumerate}
The Clay Millennium Prize mass-gap conjecture is therefore
\textbf{not proved} in this appendix. It is reduced to a single
genuinely open problem: the sharp gauge-orbit isoperimetric
inequality (the boundary case $\epsilon=1/2$ of
Theorem~\ref{thm:sharp_per_scale_cheeger}).
\end{remark}

\begin{remark}[The shortest credible route to the mass gap]
\label{rem:app226_shortest_route}
The sharp gauge-orbit isoperimetric inequality
(Remark~\ref{rem:app226_open}) is the shortest credible route
from the current theorem boundary to a genuine mass-gap proof.
It is a problem in spectral geometry on the gauge orbit
$\mathcal{A}_{j}/\mathcal{G}_{j}$, specifically the optimal
Faddeev--Popov determinant bound. The expected answer is
$h_{j}\ge c\,2^{j}/(L_{\mathrm{phys}}\sqrt{j})$ (the boundary
case $\epsilon=1/2$), which would give the mass gap
$m_{\mathrm{phys}}\ge c^{2}/(2\mathfrak{D}_{*})>0$ via
Theorem~\ref{thm:clay_rigorous}. The verification of this bound
is left as an open problem for future work.
\end{remark}

\endinput
